\documentclass[11pt,a4paper]{article}
\usepackage[utf8]{inputenc}
\usepackage[spanish,es-noquot]{babel}
\usepackage{amsmath,amssymb}
\usepackage{geometry}
\usepackage{xcolor}
\usepackage{enumitem}
\usepackage{multicol} % <--- PAQUETE AÑADIDO PARA MULTICOLUMNAS
\usepackage{tikz}
\usetikzlibrary{patterns,arrows.meta}

% Desactivar atajos conflictivos de babel en el preámbulo
\spanishdeactivate{<>}

% Configuración de márgenes (reducidos ligeramente para aprovechar mejor el papel)
\geometry{top=1.5cm, bottom=1.5cm, left=1.5cm, right=1.5cm}

% Definición de colores
\definecolor{numcolor}{rgb}{0.1, 0.4, 0.7}
\definecolor{mainblue}{rgb}{0.1, 0.4, 0.7}
\definecolor{boxbg}{rgb}{0.93, 0.95, 0.98}

% Estilo de cajas de sección
\newcommand{\sectionbox}[1]{%
  \vspace{0.15cm}\noindent
  \fcolorbox{mainblue}{boxbg}{\parbox{\dimexpr\linewidth-2\fboxsep-2\fboxrule\relax}{\textbf{\color{mainblue}\small #1}}}%
  \vspace{0.2cm}%
}

% Ajustes de separación entre columnas y listas para ahorrar espacio
\setlength{\columnsep}{1cm}
\setlist[enumerate]{itemsep=4pt, topsep=2pt, leftmargin=*}

\begin{document}

\begin{center}
    {\LARGE \textbf{Identidades notables con raíces}}\\[0.2cm]
\end{center}

\sectionbox{
    $(A + B)^2 = A^2 + 2AB + B^2$ \hfill
    $(A - B)^2 = A^2 - 2AB + B^2$ \hfill
    $(A + B)(A - B) = A^2 - B^2$ 
}

\textbf{1. Desarrolla y simplifica las siguientes operaciones con radicales:}

\begin{multicols}{2} % <--- INICIO DE 2 COLUMNAS (cambia a {3} si quieres 3 columnas)
\begin{enumerate}[label=\textbf{\arabic*.}]
    \item $\left(2\sqrt{2}\right)^2 =$
    \item $\left(3\sqrt{5}\right)^2 =$
    \item $\left(1 + \sqrt{2}\right)^2 =$
    \item $\left(\sqrt{2} + \sqrt{3}\right)^2 =$
    \item $\left(\sqrt{3} - \sqrt{2}\right)^2 =$
    \item $\left(\sqrt{2} + 1\right)\left(\sqrt{2} - 1\right) =$
    \item $\left(\sqrt{3} + \sqrt{2}\right)\left(\sqrt{3} - \sqrt{2}\right) =$
    \item $\left(1 + \sqrt{2}\right)\left(1 - \sqrt{8}\right) =$
    \item $\left(2 - \sqrt{3}\right)\left(1 + \sqrt{12}\right) =$
    \item $2\sqrt{3} \cdot 3\sqrt{2} =$
    \item $2\sqrt{8} \cdot 8\sqrt{2} =$
    \item $3\sqrt{6} \cdot 2\sqrt{3} =$
    \item $2\sqrt{15} \cdot 3\sqrt{15} =$
    \item $\left(5\sqrt{3}\right)^2 =$
    \item $\left(5 + \sqrt{3}\right)^2 =$
    \item $\left(5 - \sqrt{3}\right)^2 =$
    \item $\left(5 + \sqrt{3}\right)\left(5 - \sqrt{3}\right) =$
    \item $\left(\sqrt{5} + \sqrt{3}\right)^2 =$
    \item $\left(\sqrt{5} - \sqrt{3}\right)^2 =$
    \item $\left(3\sqrt{2} + 2\sqrt{3}\right)^2 =$
    \item $\left(2\sqrt{3} + 3\sqrt{2}\right)\left(2\sqrt{3} - 3\sqrt{2}\right) =$
    \item $\sqrt{2}\left(\sqrt{2} - 4\right) =$
    \item $\left(2 - \sqrt{3}\right)\sqrt{3} =$
    \item $\left(3\sqrt{2} + 2\right)\left(2\sqrt{3} - \sqrt{6}\right) =$
    \item $\left(2\sqrt{5} - 5\right)\sqrt{5} =$
    \item $\left(\sqrt{2} - 3\sqrt{3}\right)\left(\sqrt{2} + 5\sqrt{3}\right) =$
    \item $\left(2\sqrt{8} + 3\sqrt{2}\right)\left(3\sqrt{8} - 2\sqrt{2}\right) =$
    \item $\left(2\sqrt{5} - 5\sqrt{2}\right)\left(2\sqrt{5} + 5\sqrt{2}\right) =$
    \item $\left(2\sqrt{5} - 5\sqrt{2}\right)\left(3\sqrt{2} + 2\right) =$
    \item $\left(2\sqrt{27} - 3\right)\left(1 + \sqrt{3}\right) =$
    \item $\left(3\sqrt{8} - 4\sqrt{2}\right)\left(2\sqrt{2} - 5\sqrt{8}\right) =$
    \item $\left(\sqrt{6} + \sqrt{5}\right)^2 + \left(\sqrt{6} - \sqrt{5}\right)^2 =$
    \item $\left(\sqrt{7} + \sqrt{3}\right)^2 \left(5 - \sqrt{21}\right) =$
    \item $\left(3\sqrt{8} + 2\sqrt{2}\right)\left(2\sqrt{8} - 3\sqrt{2}\right) =$
    \item $\left(2\sqrt{3} - 3\sqrt{2}\right)^2 =$
    \item $\left(\sqrt{2} + \sqrt{3} - \sqrt{5}\right)\left(\sqrt{3} - \sqrt{2}\right) =$
\end{enumerate}
\end{multicols} % <--- FIN DE COLUMNAS


\textbf{2. Desarrolla y simplifica las siguientes expresiones aplicando las identidades notables:}

\begin{multicols}{2}
\begin{enumerate}[label=\textbf{\alph*)}]
    % Nivel básico: binomios al cuadrado y suma por diferencia con x
    \item $\left(x + 3\right)^2 =$
    \item $\left(2x - 5\right)^2 =$
    \item $\left(3x + 4\right)\left(3x - 4\right) =$
    \item $\left(x^2 + 2\right)^2 =$
    \item $\left(4x^2 - 3x\right)^2 =$
    \item $\left(5x^3 + 2x\right)\left(5x^3 - 2x\right) =$
    
    % Nivel intermedio: coeficientes fraccionarios
    \item $\left(\frac{x}{2} + 3\right)^2 =$
    \item $\left(2x - \frac{1}{3}\right)^2 =$
    \item $\left(\frac{3x}{4} + \frac{2}{5}\right)\left(\frac{3x}{4} - \frac{2}{5}\right) =$
    
    % Nivel intermedio-alto: identidades notables combinando x y radicales
    \item $\left(x + \sqrt{2}\right)^2 =$
    \item $\left(2x - \sqrt{3}\right)^2 =$
    \item $\left(x + \sqrt{5}\right)\left(x - \sqrt{5}\right) =$
    \item $\left(3x + 2\sqrt{2}\right)\left(3x - 2\sqrt{2}\right) =$
    \item $\left(x\sqrt{2} + \sqrt{3}\right)^2 =$
    \item $\left(2x\sqrt{3} - \sqrt{5}\right)^2 =$
    \item $\left(\sqrt{x} + 4\right)^2 =$ \quad $(x \ge 0)$
    \item $\left(\sqrt{2x} - \sqrt{3}\right)\left(\sqrt{2x} + \sqrt{3}\right) =$ \quad $(x \ge 0)$
    
    % Nivel avanzado: operaciones combinadas
    \item $\left(x + \sqrt{3}\right)^2 - \left(x - \sqrt{3}\right)^2 =$
    \item $\left(2x + \sqrt{5}\right)\left(2x - \sqrt{5}\right) + \left(x - \sqrt{5}\right)^2 =$
    \item $\left(x\sqrt{2} + 1\right)^2 - \left(x\sqrt{2} + 3\right)\left(x\sqrt{2} - 3\right) =$
\end{enumerate}
\end{multicols}

\vspace{0.2cm}
\textbf{3. Desarrolla y simplifica las siguientes expresiones con radicales algebraicos (considerando radicandos $\ge 0$):}

\begin{multicols}{2}
\begin{enumerate}[label=\textbf{\alph*)}]
    % Nivel 1: Radicandos simples con x
    \item $\left(\sqrt{x} + 3\right)^2 =$
    \item $\left(\sqrt{3x} - 2\right)^2 =$
    \item $\left(2\sqrt{x} + 5\right)\left(2\sqrt{x} - 5\right) =$
    \item $\left(\sqrt{5x} + \sqrt{2}\right)^2 =$
    \item $\left(3\sqrt{2x} - \sqrt{3}\right)^2 =$

    % Nivel 2: x dentro de ambos términos
    \item $\left(\sqrt{x} + \sqrt{2x}\right)^2 =$
    \item $\left(\sqrt{3x} - \sqrt{x}\right)^2 =$
    \item $\left(\sqrt{5x} + \sqrt{3x}\right)\left(\sqrt{5x} - \sqrt{3x}\right) =$
    \item $\left(x + \sqrt{x}\right)^2 =$
    \item $\left(x - \sqrt{2x}\right)\left(x + \sqrt{2x}\right) =$

    % Nivel 3: Radicandos compuestos (binomios dentro de la raíz)
    \item $\left(\sqrt{x + 1} + 2\right)^2 =$
    \item $\left(\sqrt{2x - 3} - 1\right)^2 =$
    \item $\left(\sqrt{x + 4} + \sqrt{x}\right)\left(\sqrt{x + 4} - \sqrt{x}\right) =$
    \item $\left(\sqrt{x + 1} + \sqrt{x - 1}\right)^2 =$
    \item $\left(\sqrt{2x + 1} - \sqrt{2x - 1}\right)^2 =$
\end{enumerate}
\end{multicols}

\end{document}